% Einheit 1 von 4 – Länge, Masse, Geld (bank/einheiten/e1.jsonl, zusammenbau v0.1)
%% TODO Zweigzeile Teil 1: was hier gelernt wird (Satz in Schülersprache aus dem Einheitstitel) – Einheit 1
\einheitenkopf[e1]{Einheit 1 von 4 · Länge, Masse, Geld}
\zweigzeile{ab Kl. 5 · P10}
\setcounter{aufgabe}{13}

%% TODO Titel: Ich-kann-Satz fehlt in der Bank (Platzhalter: Kettenname) – Nr. 14
%% TODO Anweisung: ein Satz über der ersten Teilaufgabe fehlt in der Bank – Nr. 14
\begin{aufgabe}{Passt die Angabe?}
\begin{teile}
\teil Wie hoch ist eine Zimmertür? Kreuze an.\\ \kreuz{$2$ m}\\ \kreuz{$9$ m}\\ \kreuz{$25$ cm}
\end{teile}
\end{aufgabe}

%% TODO \verfahren{…}: Verfahrensüberschrift in Sachsprache fehlt in der Bank (2.3 b)
%% TODO Titel: Ich-kann-Satz fehlt in der Bank (Platzhalter: Kettenname) – Nr. 15
%% TODO Anweisung: ein Satz über der ersten Teilaufgabe fehlt in der Bank – Nr. 15
\begin{aufgabe}{Länge, Masse und Geld umrechnen}
\begin{teile}
\teil Welche Einheit ist größer? Kreuze an.\\ \kreuz{Meter}\\ \kreuz{Millimeter}
\end{teile}
%% TODO Darstellung für schwach fehlt in der Bank (einheiten-e1-k2-s1-v1; Abschnitt „Für schwache Schüler“)
\swz{$7$ m – wie viel dm?}{}
%% TODO Darstellung für schwach fehlt in der Bank (einheiten-e1-k2-s1-v2; Abschnitt „Für schwache Schüler“)
\swz{$9$ cm – wie viel mm?}{}
%% TODO Darstellung für schwach fehlt in der Bank (einheiten-e1-k2-s1-v3; Abschnitt „Für schwache Schüler“)
\swz{$4$ dm – wie viel cm?}{}
%% TODO Darstellung für schwach fehlt in der Bank (einheiten-e1-k2-s1-v4; Abschnitt „Für schwache Schüler“)
\swz{$12$ m – wie viel dm?}{}
%% TODO Darstellung für schwach fehlt in der Bank (einheiten-e1-k2-s1-v5; Abschnitt „Für schwache Schüler“)
\swz{$15$ cm – wie viel mm?}{}
\swfrage{Was bleibt gleich, was ändert sich?}
%% TODO Darstellung für schwach fehlt in der Bank (einheiten-e1-k2-s2-v1; Abschnitt „Für schwache Schüler“)
\swz{$30$ mm – wie viel cm?}{}
%% TODO Darstellung für schwach fehlt in der Bank (einheiten-e1-k2-s3-v1; Abschnitt „Für schwache Schüler“)
\swz{$7\,500$ mm – wie viel m?}{}
%% TODO Darstellung für schwach fehlt in der Bank (einheiten-e1-k2-s4-v1; Abschnitt „Für schwache Schüler“)
\swz{$3$ km – wie viel m?}{}
%% TODO Darstellung für schwach fehlt in der Bank (einheiten-e1-k2-s5-v1; Abschnitt „Für schwache Schüler“)
\swz{$5$ t – wie viel kg?}{}
%% TODO Darstellung für schwach fehlt in der Bank (einheiten-e1-k2-s6-v1; Abschnitt „Für schwache Schüler“)
\swz{$3{,}45$ € – wie viel ct?}{}
%% TODO Darstellung für schwach fehlt in der Bank (einheiten-e1-k2-s7-v1; Abschnitt „Für schwache Schüler“)
\swz{$3$ m $45$ cm – wie viel cm?}{}
%% TODO Darstellung für schwach fehlt in der Bank (einheiten-e1-k2-s8-v1; Abschnitt „Für schwache Schüler“)
\swz{$2{,}35$ m – wie viel cm?}{}
%% TODO Darstellung für schwach fehlt in der Bank (einheiten-e1-k2-s9-v1; Abschnitt „Für schwache Schüler“)
\swz{Setze $<$, $>$ oder $=$ ein: $0{,}04$ m \leerfeld $40$ cm \hfill (P10 2019 OS)}{}
\swz[3]{Ein Fernseher hat eine Bildschirmdiagonale von $55$ Zoll; $1$ Zoll $\approx 2{,}54$ cm. Weise nach, dass die Diagonale etwa $1{,}40$ m lang ist. \hfill (P10 2016 OS)}{}
\end{aufgabe}

%% TODO \verfahren{…}: Verfahrensüberschrift in Sachsprache fehlt in der Bank (2.3 b)
%% TODO Titel: Ich-kann-Satz fehlt in der Bank (Platzhalter: Kettenname) – Nr. 16
%% TODO Anweisung: ein Satz über der ersten Teilaufgabe fehlt in der Bank – Nr. 16
\begin{aufgabe}{Maßstab im Koordinatensystem, Sek II}
\begin{teile}
\teil Eine Längeneinheit entspricht $30$ cm. Gesucht ist die Höhe eines Regals. Unterstreiche den Maßstabssatz. Was wird umgerechnet? Kreuze an.\\ \kreuz{Länge}\\ \kreuz{Fläche}\\ \kreuz{Volumen}
\end{teile}
%% TODO Darstellung für schwach fehlt in der Bank (einheiten-e1-k3-s1-v1; Abschnitt „Für schwache Schüler“)
\swz{Eine Leiste ist $6$ Längeneinheiten lang; eine Längeneinheit entspricht $25$ cm. Wirkliche Länge?}{}
%% TODO Darstellung für schwach fehlt in der Bank (einheiten-e1-k3-s1-v2; Abschnitt „Für schwache Schüler“)
\swz{Ein Rohr ist $7$ Längeneinheiten lang; eine Längeneinheit entspricht $15$ cm. Wirkliche Länge?}{}
%% TODO Darstellung für schwach fehlt in der Bank (einheiten-e1-k3-s1-v3; Abschnitt „Für schwache Schüler“)
\swz{Eine Kante ist $2{,}4$ Längeneinheiten lang; eine Längeneinheit entspricht $30$ cm. Wirkliche Länge?}{}
%% TODO Darstellung für schwach fehlt in der Bank (einheiten-e1-k3-s1-v4; Abschnitt „Für schwache Schüler“)
\swz{Ein Schlitz ist $9$ Längeneinheiten lang; eine Längeneinheit entspricht $0{,}2$ cm. Wirkliche Länge?}{}
%% TODO Darstellung für schwach fehlt in der Bank (einheiten-e1-k3-s1-v5; Abschnitt „Für schwache Schüler“)
\swz{Ein Griff ist $3{,}2$ Längeneinheiten lang; eine Längeneinheit entspricht $5$ cm. Wirkliche Länge?}{}
\swfrage{Was bleibt gleich, was ändert sich?}
%% TODO Darstellung für schwach fehlt in der Bank (einheiten-e1-k3-s2-v1; Abschnitt „Für schwache Schüler“)
\swz{Die Kante eines Bauteils läuft von $A(1|2)$ bis $B(1|7)$; eine Längeneinheit entspricht $4$ cm. Wie lang ist die Kante?}{}
%% TODO Darstellung für schwach fehlt in der Bank (einheiten-e1-k3-s3-v1; Abschnitt „Für schwache Schüler“)
\swz{Ein Ball fliegt auf der Bahn $f(x) = -0{,}25x^2 + 2x + 1{,}5$ ($x$: waagerechter Abstand vom Abwurf in m, $f(x)$: Höhe in m). Ein Zaun steht $3$ m vom Abwurf entfernt. In welcher Höhe fliegt der Ball über den Zaun?}{}
%% TODO Darstellung für schwach fehlt in der Bank (einheiten-e1-k3-s4-v1; Abschnitt „Für schwache Schüler“)
\swz{Eine Vase entsteht durch Drehung von $f(x) = 0{,}25x^2 + 1$ um die $x$-Achse; die Öffnung liegt bei $x = 2$, eine Längeneinheit entspricht $3$ cm. Durchmesser der Öffnung?}{}
%% TODO Darstellung für schwach fehlt in der Bank (einheiten-e1-k3-s5-v1; Abschnitt „Für schwache Schüler“)
\swz[3]{Der Bogen einer Toreinfahrt wird durch $f(x) = 0{,}6x^4 - 3x^2 + 2{,}4$ zwischen den Nullstellen $x = -1$ und $x = 1$ beschrieben; eine Längeneinheit entspricht $1{,}6$ m. Im Tor stehen $7$ senkrechte Stangen im Abstand von $40$ cm, die mittlere in der Tormitte. Berechne die Länge der vorletzten Stange rechts und die Länge einer waagerechten Stange in $2{,}4$ m Höhe. \hfill (FHR 2025)}{}
\end{aufgabe}

%% TODO Titel: Ich-kann-Satz fehlt in der Bank (Platzhalter: Kettenname) – Nr. 17
%% TODO Anweisung: ein Satz über der ersten Teilaufgabe fehlt in der Bank – Nr. 17
\begin{aufgabe}{Einheiten der Länge, der Masse und des Geldes der Größe nach ordnen}
%% TODO Darstellung für schwach fehlt in der Bank (einheiten-e1-k4-s1-v1; Abschnitt „Für schwache Schüler“)
\swz{Ordne, mit der kleinsten beginnend: m, mm, km, cm, dm}{}
\end{aufgabe}

%% TODO Titel: Ich-kann-Satz fehlt in der Bank (Platzhalter: Kettenname) – Nr. 18
%% TODO Anweisung: ein Satz über der ersten Teilaufgabe fehlt in der Bank – Nr. 18
\begin{aufgabe}{Repräsentanten zuordnen (welche Angabe passt zu welchem Gegenstand)}
%% TODO Darstellung für schwach fehlt in der Bank (einheiten-e1-k5-s1-v1; Abschnitt „Für schwache Schüler“)
\swz{Ordne zu: Heft, Fahrrad, Fußballfeld – $105$ m; $30$ cm; $1{,}8$ m}{}
\end{aufgabe}

%% TODO Titel: Ich-kann-Satz fehlt in der Bank (Platzhalter: Kettenname) – Nr. 19
%% TODO Anweisung: ein Satz über der ersten Teilaufgabe fehlt in der Bank – Nr. 19
\begin{aufgabe}{drei bis vier Größen ordnen}
%% TODO Darstellung für schwach fehlt in der Bank (einheiten-e1-k6-s1-v1; Abschnitt „Für schwache Schüler“)
\swz{Ordne, mit der kleinsten beginnend: $1{,}2$ m; $95$ cm; $1\,150$ mm}{}
\end{aufgabe}

%% TODO Titel: Ich-kann-Satz fehlt in der Bank (Platzhalter: Kettenname) – Nr. 20
%% TODO Anweisung: ein Satz über der ersten Teilaufgabe fehlt in der Bank – Nr. 20
\begin{aufgabe}{Umrechnen mit gegebenem Faktor bei nichtmetrischen Einheiten (Fuß, Meile, Zoll; P10-Form)}
%% TODO Darstellung für schwach fehlt in der Bank (einheiten-e1-k7-s1-v1; Abschnitt „Für schwache Schüler“)
\swz{Ein Monitor misst $32$ Zoll in der Diagonale; $1$ Zoll $\approx 2{,}54$ cm. Wie viel cm? Runde auf eine Stelle.}{}
\end{aufgabe}

%% TODO Titel: Ich-kann-Satz fehlt in der Bank (Platzhalter: Kettenname) – Nr. 21
%% TODO Anweisung: ein Satz über der ersten Teilaufgabe fehlt in der Bank – Nr. 21
\begin{aufgabe}{Größen mit gleicher Einheit addieren und subtrahieren}
%% TODO Darstellung für schwach fehlt in der Bank (einheiten-e1-k8-s1-v1; Abschnitt „Für schwache Schüler“)
\swz{$3{,}45$ m $+ 1{,}8$ m?}{}
\end{aufgabe}

%% TODO Titel: Ich-kann-Satz fehlt in der Bank (Platzhalter: Kettenname) – Nr. 22
%% TODO Anweisung: ein Satz über der ersten Teilaufgabe fehlt in der Bank – Nr. 22
\begin{aufgabe}{Länge, Masse und Geld umrechnen – Fehler finden, Begründen, Anwendung, Darstellungswechsel}
\swz[3]{Tom schreibt: $450$ cm $= 45\,000$ m. Wo ist der Fehler?}{}
\swz[3]{Eine Längeneinheit entspricht $0{,}4$ cm. Ein Bauteil ist $7$ Längeneinheiten hoch. Jan schreibt: Das Bauteil ist $7$ cm hoch. Wo ist der Fehler?}{}
\swfrage{Begründe: Warum ist $0{,}09$ m weniger als $90$ cm?}
\swfrage{Anna sagt: „Eine Längeneinheit ist immer ein Zentimeter.“ Begründe, warum das nicht stimmt.}
\swz[3]{Ein Regal soll in eine Nische von $1{,}35$ m Breite. Die Böden sind $128$ cm lang, dazu kommen zwei Seitenwände von je $2$ cm. Passt das Regal hinein?}{}
\swa{Schreibe die Länge aus der Stellenwerttafel als Kommazahl in m und als Angabe in m und cm. \leerfeld[m] ; \leerfeld[m] \leerfeld[cm]}{\sachtabelle{cccc}{m & dm & cm & mm}{2 & 0 & 6 & 0}}
\end{aufgabe}

%% TODO Merkkasten Einheit 1: 6 Zeilen, Vorgabe höchstens fünf (3.1)
\uebersichtskasten{Einheiten ordnen: Jede Einheit hat eine Nachbareinheit. Von einer großen Einheit zur kleinen nimmst du mal, von der kleinen zur großen teilst du. \\ Länge: mm – cm – dm – m (Schritt 10), m – km (Schritt 1000). 3 m = 300 cm; 4500 m = 4,5 km. \\ Masse: mg – g – kg – t (Schritt 1000). 2,5 kg = 2500 g. \\ Geld: 1 € = 100 ct. 0,58 € = 58 ct. \\ Vergleichen: Erst beide Angaben in dieselbe Einheit bringen, dann die Zahlen vergleichen. 0,06 m = 6 cm, also 0,06 m < 60 cm. \\ Maßstab im Koordinatensystem (Sek II): Der Text legt fest, was eine Längeneinheit wirklich ist – Länge in Längeneinheiten mal Faktor. Eine Längeneinheit entspricht 20 cm: eine Leiste von 4 LE ist 4 · 20 cm = 80 cm lang; eine Längeneinheit entspricht 0,5 cm: eine Höhe von 4 LE ist 2 cm.}
